Lässt man einen \mStO\ schweren Stein in einen \hBrO\ tiefen Brunnenschacht fallen, so nimmt sein Impuls während des Falls ständig zu.

\begin{abcliste}
\abc Wie lässt sich dieser scheinbare Widerspruch zum Impulserhaltungssatz auflösen?
\abc Wie gross ist der Geschwindigkeitszuwachs der Erde beim Aufprall des Steins? Wie interpretieren Sie dieses Resultat?
\end{abcliste}
